\documentclass{standalone}
\usepackage{circuitikz}
\usetikzlibrary{calc}
\definecolor{circuitnotemark}{HTML}{FE00FE}
\begin{document}
\begin{circuitikz}[american, line width=0.8pt]
\ctikzset{bipoles/length=1.2cm}
\ctikzset{grounds/scale=1.36}
\foreach \x in {0,2,4} {\foreach \y in {0,-2} {\fill[gray, opacity=0.35] (\x,\y) circle (1.1pt);}}
\node[gray, font=\scriptsize] at (0,0.84) {1};
\node[gray, font=\scriptsize] at (2,0.84) {2};
\node[gray, font=\scriptsize] at (4,0.84) {3};
\node[gray, font=\scriptsize, anchor=east] at (-0.84,0) {a};
\node[gray, font=\scriptsize, anchor=east] at (-0.84,-2) {b};
\coordinate (a1) at (0,0);
\coordinate (a2) at (2,0);
\coordinate (b2) at (2,-2);
\coordinate (a3) at (4,0);
\draw (a1) node[ocirc]{} node[above left]{IN}; % line 3
\draw (a1) to[R, l_=$R_{1}$, a^=$10\,\mathrm{k}\Omega$] (a2); % line 4
\draw (a2) to[C, l_=$C_{1}$, a^=$100\,\mathrm{n}\mathrm{F}$] (b2); % line 5
\draw (a3) node[ocirc]{} node[above left]{OUT}; % line 6
\node[ground] at (b2) {}; % line 7
\draw (a2) -- (a3); % line 9
\node[circ] at (a2) {};
\path (6,-5.263) rectangle (10.724,0.395); % line 11
\node[anchor=west, inner sep=0, circuitnotemark, font=\footnotesize] at (6,0) {X}; % line 11
\node[anchor=west, inner sep=0, circuitnotemark, font=\footnotesize] at (6,-0.325) {X}; % line 11
\node[anchor=west, inner sep=0, circuitnotemark, font=\footnotesize] at (6,-0.649) {X}; % line 11
\node[anchor=west, inner sep=0, circuitnotemark, font=\footnotesize] at (6,-0.974) {X}; % line 11
\node[anchor=west, inner sep=0, circuitnotemark, font=\footnotesize] at (6,-1.298) {X}; % line 11
\node[anchor=west, inner sep=0, circuitnotemark, font=\footnotesize] at (6,-1.623) {X}; % line 11
\node[anchor=west, inner sep=0, circuitnotemark, font=\footnotesize] at (6,-1.947) {X}; % line 11
\node[anchor=west, inner sep=0, circuitnotemark, font=\footnotesize] at (6,-2.272) {X}; % line 11
\node[anchor=west, inner sep=0, circuitnotemark, font=\footnotesize] at (6,-2.596) {X}; % line 11
\node[anchor=west, inner sep=0, circuitnotemark, font=\footnotesize] at (6,-2.921) {X}; % line 11
\node[anchor=west, inner sep=0, circuitnotemark, font=\footnotesize] at (6,-3.246) {X}; % line 11
\node[anchor=west, inner sep=0, circuitnotemark, font=\footnotesize] at (6,-3.57) {X}; % line 11
\node[anchor=west, inner sep=0, circuitnotemark, font=\footnotesize] at (6,-3.895) {X}; % line 11
\node[anchor=west, inner sep=0, circuitnotemark, font=\footnotesize] at (6,-4.219) {X}; % line 11
\node[anchor=west, inner sep=0, circuitnotemark, font=\footnotesize] at (6,-4.544) {X}; % line 11
\node[anchor=west, inner sep=0, circuitnotemark, font=\footnotesize] at (6,-4.868) {X}; % line 11
\node[anchor=south west, inner sep=2pt, circuitnotemark, font=\large\bfseries] (circuittitle) at (current bounding box.north west) {X};
\path (circuittitle.south west) rectangle ++(3.429,0.593);
\node[gray, anchor=north east, font=\tiny, inner sep=2pt] at (current bounding box.south east) {circuit-fence 0.1.0};
\end{circuitikz}
\end{document}
